\DocumentMetadata{lang=en-US,pdfstandard=A-2b,uncompress}
\documentclass[12pt]{article}
\usepackage[utf8]{inputenc}
\usepackage[margin=1in]{geometry}
\usepackage{amsmath,amssymb}
\usepackage{enumitem}
\usepackage{parskip}
\usepackage{hyperref}
\usepackage[tagged, highstructure]{accessibility}

\title{Quiz 2 Answer Key: DNF and Soundness}
\hypersetup{pdftitle={Quiz 2 Answer Key: DNF and Soundness}}
\author{}
\date{}

\newcommand{\forces}{\models}
\newcommand{\meta}[1]{\mathcal{#1}}

\begin{document}
\maketitle

\begin{enumerate}
    \item
    \begin{enumerate}[label=(\alph*)]
        \item For
        \[
        (P \rightarrow R)\leftrightarrow R
        \]
        the truth table is:
        \[
        \begin{array}{c c | c | c}
        P & R & (P\rightarrow R) & (P\rightarrow R)\leftrightarrow R\\
        \hline
        T & T & T & T\\
        T & F & F & T\\
        F & T & T & T\\
        F & F & T & F
        \end{array}
        \]
        True rows are \((T,T)\), \((T,F)\), \((F,T)\). So one DNF is:
        \[
        (P\land R)\vee(P\land\neg R)\vee(\neg P\land R).
        \]

        \item Why this works:
        each disjunct corresponds to one row where the original sentence is true, and is true on exactly that row. Disjoining them gives a sentence true on exactly the same rows as the original sentence. Number of disjuncts = number of true rows.
    \end{enumerate}

    \item
    \begin{enumerate}[label=(\alph*)]
        \item
        \[
        \phi_6=\neg P,\quad \delta_6=\{P\rightarrow Q,\ \neg Q\}.
        \]
        \[
        \phi_8=(P\rightarrow Q)\rightarrow(\neg Q\rightarrow\neg P),\quad \delta_8=\varnothing.
        \]
        (No premises are left in force on line 8.)

        \item Line \(n\) is shiny iff \(\delta_n\forces \phi_n\), i.e. every valuation satisfying all assumptions in force at line \(n\) satisfies the line-\(n\) sentence.

        \item Rule \(R\) is rule-sound iff whenever all lines \(<n\) are shiny and line \(n\) is justified by \(R\), then line \(n\) is shiny.
    \end{enumerate}

    \item Fill-in targets:
    \begin{enumerate}[label=\arabic*.]
        \item By rule structure + closed-subproof constraint: the assumptions in force at \(j\) are \(\neg A\) plus assumptions still available at \(n\), so \(\exists\delta'\subseteq\delta_n\) with \(\delta_j=\delta'\cup\{\neg A\}\).
        \item By IH (since \(j<n\)).
        \item By truth conditions for negation: if \(v\not\forces A\), then \(v\forces\neg A\).
        \item From \(v\forces\delta_n\), \(\delta'\subseteq\delta_n\), and \(v\forces\neg A\), we get \(v\forces\delta'\cup\{\neg A\}=\delta_j\).
        \item From \(\delta_j\forces\bot\) and \(v\forces\delta_j\), it follows that \(v\forces\bot\), contradiction.
    \end{enumerate}

    \item \( \vee I \) is rule-sound (sample proof).

    Assume all lines \(<n\) are shiny and line \(n\) is justified by \(\vee I\) from line \(i\). WLOG:
    \[
    \phi_i=A,\qquad \phi_n=A\vee B.
    \]
    Since a cited line cannot be inside a closed subproof, \(\delta_i\subseteq\delta_n\).

    By IH, line \(i\) is shiny, so \(\delta_i\forces A\). Hence \(\delta_n\forces A\).
    Let \(v\) be arbitrary with \(v\forces\delta_n\). Then \(v\forces A\), so by the truth table for \(\vee\), \(v\forces A\vee B\).
    Therefore every \(v\) satisfying \(\delta_n\) satisfies \(\phi_n=A\vee B\), so \(\delta_n\forces\phi_n\). Thus line \(n\) is shiny, and \(\vee I\) is rule-sound.

\end{enumerate}

\end{document}
